% Options for packages loaded elsewhere
\PassOptionsToPackage{unicode}{hyperref}
\PassOptionsToPackage{hyphens}{url}
\documentclass[
  12pt,
]{ctexart}
\usepackage{xcolor}
\usepackage{amsmath,amssymb}
\setcounter{secnumdepth}{-\maxdimen} % remove section numbering
\usepackage{iftex}
\ifPDFTeX
  \usepackage[T1]{fontenc}
  \usepackage[utf8]{inputenc}
  \usepackage{textcomp} % provide euro and other symbols
\else % if luatex or xetex
  \usepackage{unicode-math} % this also loads fontspec
  \defaultfontfeatures{Scale=MatchLowercase}
  \defaultfontfeatures[\rmfamily]{Ligatures=TeX,Scale=1}
\fi
\usepackage{lmodern}
\ifPDFTeX\else
  % xetex/luatex font selection
\fi
% Use upquote if available, for straight quotes in verbatim environments
\IfFileExists{upquote.sty}{\usepackage{upquote}}{}
\IfFileExists{microtype.sty}{% use microtype if available
  \usepackage[]{microtype}
  \UseMicrotypeSet[protrusion]{basicmath} % disable protrusion for tt fonts
}{}
\makeatletter
\@ifundefined{KOMAClassName}{% if non-KOMA class
  \IfFileExists{parskip.sty}{%
    \usepackage{parskip}
  }{% else
    \setlength{\parindent}{0pt}
    \setlength{\parskip}{6pt plus 2pt minus 1pt}}
}{% if KOMA class
  \KOMAoptions{parskip=half}}
\makeatother
\usepackage{longtable,booktabs,array}
\usepackage{calc} % for calculating minipage widths
% Correct order of tables after \paragraph or \subparagraph
\usepackage{etoolbox}
\makeatletter
\patchcmd\longtable{\par}{\if@noskipsec\mbox{}\fi\par}{}{}
\makeatother
% Allow footnotes in longtable head/foot
\IfFileExists{footnotehyper.sty}{\usepackage{footnotehyper}}{\usepackage{footnote}}
\makesavenoteenv{longtable}
\setlength{\emergencystretch}{3em} % prevent overfull lines
\providecommand{\tightlist}{%
  \setlength{\itemsep}{0pt}\setlength{\parskip}{0pt}}
\usepackage{bookmark}
\IfFileExists{xurl.sty}{\usepackage{xurl}}{} % add URL line breaks if available
\urlstyle{same}
\hypersetup{
  hidelinks,
  pdfcreator={LaTeX via pandoc}}

\author{SCX}
\date{2026}

\begin{document}

\section{SCX 开发 TODO ---
已归档}\label{scx-ux5f00ux53d1-todo-ux5df2ux5f52ux6863}

\begin{quote}
\textbf{本文件已被 \texttt{SCX\_TODO\_2026-06-26.md} 替代} (2026-06-26)

本文档保留仅用于历史参考，不再作为当前 TODO。后续更新请查看
\texttt{SCX\_TODO\_2026-06-26.md}。
\end{quote}

\begin{quote}
最后更新：2026-06-26 \textbar{} 基于 6 Agent 讨论的综合方案
\end{quote}

\begin{center}\rule{0.5\linewidth}{0.5pt}\end{center}

\subsection{Phase 1：理论加固（1-2
月）}\label{phase-1ux7406ux8bbaux52a0ux56fa1-2-ux6708}

\subsubsection{1.1 SCX-Compress 定理 🔴
P0}\label{scx-compress-ux5b9aux7406-p0}

\begin{itemize}
\tightlist
\item[$\boxtimes$]
  \textbf{Theorem 1} ✅ ---
  \texttt{theory/propositions/04\_compression\_fidelity.md} (689行)

  \begin{itemize}
  \tightlist
  \item
    压缩保真上界: 偏差项 + Rademacher 方差项
  \item
    核心结论: 50\% 压缩且精度损失 \textless5\% 需要 D(s) ≥ 0.90
  \end{itemize}
\item[$\square$]
  合成验证: 在 2D 合成数据上数值验证 bound 的 tightness
\item[$\square$]
  \textbf{写入 Paper 4 §4.3}
\end{itemize}

\subsubsection{1.2 Expert Governance Protocol 形式化 🔴
P0}\label{expert-governance-protocol-ux5f62ux5f0fux5316-p0}

\begin{itemize}
\tightlist
\item[$\boxtimes$]
  \textbf{5 步协议} ✅ ---
  \texttt{theory/propositions/05\_expert\_governance\_protocol.md}

  \begin{itemize}
  \tightlist
  \item
    Gauge Check → Domain Certificate → Conflict Resolution → Anchor
    Verification → Distillation Authorization
  \end{itemize}
\item[$\boxtimes$]
  \textbf{Certification Theorem} ✅ --- R\_student(s) ≤ R\_min(s) +
  ε(N\_anchor, δ)
\item[$\square$]
  在 AlN v3 数据上 retrospective 验证
\item[$\square$]
  \textbf{写入 Paper 3 §3}
\end{itemize}

\subsubsection{1.3 State Discovery 鲁棒性 🟡
P1}\label{state-discovery-ux9c81ux68d2ux6027-p1}

\begin{itemize}
\tightlist
\item[$\boxtimes$]
  \textbf{Mis-specification 分析} ✅ ---
  \texttt{src/scx/state/robustness.py} (merge/split/noise impact)
\item[$\boxtimes$]
  \textbf{跨描述符稳定性} ✅ --- \texttt{cross\_descriptor\_stability()}
  (ARI)
\item[$\boxtimes$]
  \textbf{Bootstrap 稳健 K} ✅ --- \texttt{suggest\_robust\_states()}
\item[$\boxtimes$]
  \textbf{边界置信度} ✅ --- \texttt{state\_boundary\_confidence()}
\item[$\square$]
  稳定性定理: state partition 小扰动 → 分类结果变化 bounded
\item[$\square$]
  \textbf{写入 Paper 4 §3.4}
\end{itemize}

\subsubsection{1.4 自适应阈值 🟢
P2}\label{ux81eaux9002ux5e94ux9608ux503c-p2}

\begin{itemize}
\tightlist
\item[$\boxtimes$]
  \textbf{网格搜索校准} ✅ --- \texttt{src/scx/valuation/adaptive.py}
  (\texttt{calibrate()})
\item[$\boxtimes$]
  \textbf{Sensitivity analysis} ✅ --- 阈值 ±20\% → 分类变化
\item[$\boxtimes$]
  \textbf{无监督自适应} ✅ --- percentile/gap 方法
  (\texttt{auto\_threshold()})
\item[$\square$]
  \textbf{写入 Paper 4 §5.1}
\end{itemize}

\begin{center}\rule{0.5\linewidth}{0.5pt}\end{center}

\subsection{Phase 2：实验补强（2-3
月）}\label{phase-2ux5b9eux9a8cux8865ux5f3a2-3-ux6708}

\subsubsection{2.1 SCX + Influence 融合 🟡
P1}\label{scx-influence-ux878dux5408-p1}

\begin{itemize}
\tightlist
\item[$\square$]
  实现 \texttt{scx/valuation/influence.py} --- State-Conditioned
  Influence
\item[$\square$]
  CIFAR-100 上对比: SCX vs LESS vs SCX+LESS
\item[$\square$]
  证明: SCX 粗选状态 + Influence 细选样本 \textgreater{} 任一单独方法
\item[$\square$]
  \textbf{写入 Paper 4 §6.3}
\end{itemize}

\subsubsection{2.2 通用 ML 实验 🟡
P1}\label{ux901aux7528-ml-ux5b9eux9a8c-p1}

\begin{itemize}
\tightlist
\item[$\square$]
  \textbf{CIFAR-100 + 5 ResNet-18}: 每个在不同子集训练 → multi-expert

  \begin{itemize}
  \tightlist
  \item[$\square$]
    注入 label noise (10-30\%)
  \item[$\square$]
    SCX 四分类 vs ground truth noise labels
  \item[$\square$]
    对比 Data Shapley (pyDVL), LOO, Random
  \end{itemize}
\item[$\square$]
  \textbf{UCI Tabular × 2}: 10+ 维数据集

  \begin{itemize}
  \tightlist
  \item[$\square$]
    对比 Uncertainty, Diversity, Coreset
  \end{itemize}
\item[$\square$]
  \textbf{产出 3 张 Figure} (acquisition efficiency, classification
  accuracy, ablation)
\end{itemize}

\subsubsection{2.3 自适应阈值 🟢
P2}\label{ux81eaux9002ux5e94ux9608ux503c-p2-1}

\begin{itemize}
\tightlist
\item[$\square$]
  实现 \texttt{DataClassifier.auto\_calibrate(X,\ y\_labeled)}
\item[$\square$]
  Sensitivity analysis: 阈值 ±20\% → 分类变化
\item[$\square$]
  Bayesian 阈值 (可选)
\item[$\square$]
  \textbf{写入 Paper 4 §5.1}
\end{itemize}

\begin{center}\rule{0.5\linewidth}{0.5pt}\end{center}

\subsection{Phase 3：扩展（3-6
月）}\label{phase-3ux6269ux5c553-6-ux6708}

\subsubsection{3.1 Online SCX 🔵 P3}\label{online-scx-p3}

\begin{itemize}
\tightlist
\item[$\square$]
  设计 \texttt{OnlineSCXFramework} API
\item[$\square$]
  增量状态更新 (exponential moving average)
\item[$\square$]
  在线专家可靠性追踪
\item[$\square$]
  合成数据流实验
\end{itemize}

\subsubsection{3.2 Benchmark Suite 🔵 P3}\label{benchmark-suite-p3}

\begin{itemize}
\tightlist
\item[$\square$]
  5-10 数据集 (合成 + UCI + 图像 + MLIP)
\item[$\square$]
  标准化 API: \texttt{benchmark.run(method,\ dataset)\ →\ results}
\item[$\square$]
  Baseline 集成: Random, Uncertainty, Diversity, Shapley, LESS, Coreset,
  SCX
\item[$\square$]
  Leaderboard + 文档
\end{itemize}

\begin{center}\rule{0.5\linewidth}{0.5pt}\end{center}

\subsection{Phase 4：论文写作}\label{phase-4ux8bbaux6587ux5199ux4f5c}

\subsubsection{Paper 3 (Expert Compiler) --- 等 AlN v4
重训后}\label{paper-3-expert-compiler-ux7b49-aln-v4-ux91cdux8badux540e}

\begin{itemize}
\tightlist
\item[$\square$]
  §3 Expert Compilation Protocol (方向 4)
\item[$\square$]
  §5 Experiments: retrospective on AlN v3
\end{itemize}

\subsubsection{Paper 4 (SCX) ---
持续更新}\label{paper-4-scx-ux6301ux7eedux66f4ux65b0}

\begin{itemize}
\tightlist
\item[$\square$]
  §3.4 Robustness Analysis (方向 2)
\item[$\square$]
  §4.3 Compression Theory (方向 1)
\item[$\square$]
  §5.1 Adaptive Threshold (方向 6)
\item[$\square$]
  §6.2 General ML Experiments (方向 5)
\item[$\square$]
  §6.3 Comparison with Influence Methods (方向 3)
\end{itemize}

\begin{center}\rule{0.5\linewidth}{0.5pt}\end{center}

\subsection{快速索引}\label{ux5febux901fux7d22ux5f15}

\begin{longtable}[]{@{}
  >{\raggedright\arraybackslash}p{(\linewidth - 2\tabcolsep) * \real{0.4500}}
  >{\raggedright\arraybackslash}p{(\linewidth - 2\tabcolsep) * \real{0.5500}}@{}}
\toprule\noalign{}
\begin{minipage}[b]{\linewidth}\raggedright
想做什么
\end{minipage} & \begin{minipage}[b]{\linewidth}\raggedright
看哪个文件
\end{minipage} \\
\midrule\noalign{}
\endhead
\bottomrule\noalign{}
\endlastfoot
完整扩展方案 & \texttt{SCX\_思想扩展\_综合方案.md} \\
核心数学 & \texttt{agent\_outputs/01\_SCX\_核心框架\_数学分析.md} \\
子模块设计 & \texttt{agent\_outputs/02\_SCX\_子模块\_详细设计.md} \\
竞争分析: 数据估值 &
\texttt{agent\_outputs/03\_竞争分析\_数据估值与主动学习.md} \\
竞争分析: MoE/蒸馏 &
\texttt{agent\_outputs/04\_竞争分析\_MoE与蒸馏.md} \\
数学根源 & \texttt{agent\_outputs/05\_数学根源与证明.md} \\
实现架构 & \texttt{agent\_outputs/06\_实现架构与实验设计.md} \\
Python 代码 & \texttt{../src/scx/} \\
测试 & \texttt{../tests/} \\
\end{longtable}

\end{document}
